\documentclass[12pt,a4paper]{article}
\usepackage[utf-8]{inputenc}
\usepackage[T1]{fontenc}
\usepackage[french]{babel}
\usepackage[margin=2.5cm]{geometry}
\usepackage{hyperref}
\usepackage{graphicx}
\usepackage{pdfpages}
\usepackage{tocloft}
\usepackage{fancyhdr}
\usepackage{setspace}

% Configuration des couleurs et du style
\usepackage{xcolor}
\definecolor{darkblue}{rgb}{0.1,0.2,0.5}
\definecolor{lightgray}{rgb}{0.95,0.95,0.95}

% En-têtes et pieds de page
\pagestyle{fancy}
\fancyhf{}
\rhead{\textit{PFE - Flood Damage XAI MLOps}}
\lfoot{\today}
\rfoot{\thepage}
\renewcommand{\headrulewidth}{0.5pt}

% Titre principal
\title{
    \textbf{\Large Classification spatio-temporelle des dégâts d'inondation} \\
    \textbf{\Large à partir d'images satellites avec XAI et MLOps}
    \\[0.5cm]
    \normalsize Projet de Fin d'Études (PFE)
}
\author{}
\date{\today}

\begin{document}

% Page de titre
\maketitle
\thispagestyle{empty}

\vspace{2cm}

\begin{center}
    \textcolor{darkblue}{\textbf{Sommaire}}
\end{center}

\vspace{1cm}

\tableofcontents

\newpage

% Section 1 : Introduction et Contexte
\section{Introduction et Contexte du Projet}

\subsection{Contexte général}
Les catastrophes naturelles, en particulier les inondations, constituent une menace majeure pour les sociétés modernes. Selon les données du Global Insure Forum, les inondations représentent environ 40\% des pertes économiques liées aux catastrophes naturelles worldwide. L'évaluation rapide et précise des dégâts post-désastre est essentielle pour les opérations de secours, la planification de la reconstruction et l'évaluation des assurances.

Les images satellites haute résolution offrent une couverture géographique complète et sont accessibles peu après un événement. Cependant, l'analyse manuelle de ces images est\:
\begin{itemize}
    \item Chronophage et coûteuse
    \item Sujette à la variabilité humaine
    \item Difficile à mettre à l'échelle
\end{itemize}

L'utilisation de modèles de deep learning offre une solution prometteuse pour automatiser et standardiser cette analyse.

\subsection{Motivation scientifique}
Ce projet s'inscrit dans la convergence de trois domaines clés\:

\subsubsection{1. Computer Vision multi-temporelle}
La classification des dégâts exige l'analyse de deux images satellites\:
\begin{itemize}
    \item \textbf{Images pré-désastre} : état de référence du bâti
    \item \textbf{Images post-désastre} : identification des modifications
\end{itemize}

Cette approche bi-temporelle est fundamentale pour distinguer les dégâts des variations naturelles (ombres, saisonnalité, etc.).

\subsubsection{2. Explicabilité et validation scientifique}
Les modèles black-box ne sont pas acceptés dans un contexte d'aide à la décision pour les catastrophes. La projet intègre Grad-CAM pour visualiser les zones que le modèle utilise pour sa classification. Cela permet de\:
\begin{itemize}
    \item Vérifier que le modèle regarde effectivement le bâti
    \item Identifier les artefacts spurieux
    \item Renforcer la confiance des stakeholders
\end{itemize}

\subsubsection{3. Machine Learning Operations (MLOps)}
Un modèle scientifiquement robuste ne suffisit pas. L'industrialisation require\:
\begin{itemize}
    \item Versionnement des données et modèles
    \item Reproductibilité des expériences
    \item Traçabilité du pipeline
    \item Déploiement et monitoring scalables
\end{itemize}

Ce projet démontre une approche intégrée du ML moderne.

\subsection{Objectifs du projet}
\subsubsection{Objectif principal}
Développer et valider une pipeline complète de classification spatio-temporelle des dégâts d'inondation, combinant\:
\begin{enumerate}
    \item Une architecture de deep learning robuste (Siamese Network)
    \item Une fondation data solide et anti-leakage
    \item L'explicabilité scientifique (XAI)
    \item L'infrastructure MLOps (versionning, reproductibilité)
\end{enumerate}

\subsubsection{Objectifs secondaires}
\begin{itemize}
    \item Démontrer l'efficacité de l'apprentissage de représentations bi-temporelles
    \item Évaluer la performance sur des classes fortement déséquilibrées
    \item Comparer différentes architectures (ResNet vs Transformer)
    \item Produire un cadre reproductible pour futurs travaux
\end{itemize}

\subsection{Architecture générale}
Le projet s'appuie sur les composants suivants\:

\subsubsection{Data Foundation}
\begin{itemize}
    \item Dataset source\: xBD (subset flooding)
    \item Métadonnées\: 8471 exemples bâtiments, 243 scènes
    \item Anti-leakage\: split au niveau scène
    \item Bi-temporalité\: synchronisation pré/post
\end{itemize}

\subsubsection{Modèles}
\begin{itemize}
    \item Architecture primaire\: Siamese ResNet18
    \item Variantes expérimentées\: ResNet18 avec régularisation, BitTemporalTransformer
    \item Gestion du déséquilibre\: WeightedRandomSampler + class weights
\end{itemize}

\subsubsection{Infrastructure MLOps}
\begin{itemize}
    \item Versionnement données\: DVC (Data Version Control)
    \item Versionnement modèles\: MLflow Model Registry
    \item Tracking d'expériences\: MLflow
    \item Validation data\: scripts d'audit complets
\end{itemize}

\subsubsection{Explicabilité}
\begin{itemize}
    \item Méthode\: Grad-CAM (Gradient-weighted Class Activation Maps)
    \item Validation qualitative\: visualisation de zones d'intérêt
    \item Audit du modèle\: détection d'artefacts spurieux
\end{itemize}

\subsubsection{Déploiement (futur)}
\begin{itemize}
    \item API\: FastAPI avec endpoints REST
    \item Interface utilisateur\: Streamlit
    \item Monitoring\: Evidently AI
\end{itemize}

\subsection{Contribution attendue}
Ce projet contribue à trois niveaux\:

\subsubsection{Au niveau scientifique}
\begin{itemize}
    \item Démonstration de l'efficacité de la synchronisation bi-temporelle
    \item Analyse de l'impact de l'architectures sur des données hautement déséquilibrées
    \item Validation de Grad-CAM comme outil de confiance scientifique
\end{itemize}

\subsubsection{Au niveau méthodologique}
\begin{itemize}
    \item Pattern de Data Foundation reproductible
    \item Bonnes pratiques d'anti-leakage pour des données temporelles
    \item Intégration MLOps pour l'analyse d'imagerie satellite
\end{itemize}

\subsubsection{Au niveau pratique}
\begin{itemize}
    \item Codebase production-ready
    \item Documentation technique complète
    \item Pipeline pré-déploiement
\end{itemize}

\subsection{Structure du mémoire}
Ce document est organisé comme suit\:

\begin{description}
    \item[Section 2] \textit{Data Foundation} : construction des métadonnées, gestion de l'anti-leakage, structure bi-temporelle
    \item[Section 3] \textit{Pipeline d'apprentissage} : architecture Siamese, stratégies de régularisation, résultats expérimentaux
    \item[Section 4] \textit{Explicabilité (XAI)} : Grad-CAM, validation qualitative, audit du modèle
    \item[Section 5] \textit{MLOps et versionnement} : DVC, MLflow, infrastructure de reproductibilité
    \item[Section 6] \textit{Résultats et interprétation} : matrices de confusion, leçons d'ingénierie
    \item[Section 7] \textit{Perspectives futures} : déploiement, monitoring, extensibilité
\end{description}

\newpage

% Section 2 : Data Foundation
\section{Data Foundation}

La data foundation constitue le socle du projet. Elle garantit :
\begin{itemize}
    \item \textbf{Reproductibilité} : métadonnées versionnées et splits robustes
    \item \textbf{Anti-leakage} : séparation scène-level (sample\_id)
    \item \textbf{Traçabilité} : chemins relatifs et validation exhaustive
    \item \textbf{Bi-temporalité} : synchronisation pré/post images
\end{itemize}

\subsection{Structure des métadonnées}
Les fichiers centraux sont :
\begin{itemize}
    \item \texttt{data/interim/metadata.csv} : 8471 bâtiments, 243 scènes
    \item \texttt{data/splits/metadata\_splits.csv} : train/val/test sans chevauchement
\end{itemize}

\subsection{Gestion du déséquilibre}
\begin{center}
    \begin{tabular}{|c|c|}
        \hline
        \textbf{Classe} & \textbf{Nombre d'exemples} \\
        \hline
        no-damage & 8128 (96.0\%) \\
        destroyed & 75 (0.9\%) \\
        major-damage & 119 (1.4\%) \\
        minor-damage & 149 (1.8\%) \\
        \hline
    \end{tabular}
\end{center}

Stratégie : \texttt{WeightedRandomSampler} avec calcul de \texttt{class\_weights}.

\newpage

% Section 3 : Pipeline d'apprentissage
\section{Pipeline d'apprentissage}

\subsection{Architecture Siamese ResNet18}
\begin{itemize}
    \item Deux branches pré/post indépendantes
    \item Poids partagés
    \item Couche d'fusion pour classification
\end{itemize}

\subsection{Runs d'expérimentation}
\begin{enumerate}
    \item \textbf{Run A} : Baseline ResNet18, régularisation L2
    \item \textbf{Run B} : RegularizedNet avec dropout et batch norm avancée
    \item \textbf{Run C} : BitTemporalTransformer (exploratoire)
\end{enumerate}

\newpage

% Section 4 : Explicabilité (XAI)
\section{Explicabilité (XAI)}

\subsection{Grad-CAM}
Grad-CAM génère des heatmaps pour visualiser les régions que le modèle utilise pour sa prédiction.

\subsection{Validation scientifique}
Les visualisations Grad-CAM permettent de vérifier que :
\begin{itemize}
    \item Le modèle regarde effectivement la zone du bâtiment
    \item Les cues de dégât post-désastre sont détectées
    \item Aucun artefact spurieux n'influence les prédictions
\end{itemize}

\newpage

% Section 5 : MLOps et versionnement
\section{MLOps et versionnement}

\subsection{DVC (Data Version Control)}
Le projet utilise DVC pour versionner explicitement les artefacts produits par la data foundation, en complément de Git qui versionne le code et les configurations. À ce stade, la mise en œuvre DVC couvre les deux tables centrales suivantes :
\begin{itemize}
    \item \texttt{data/interim/metadata.csv} : table principale des bâtiments et des trajectoires pré/post
    \item \texttt{data/splits/metadata\_splits.csv} : table des partitions \textit{train}, \textit{val}, \textit{test}
\end{itemize}
Ces fichiers sont suivis via des fichiers DVC dédiés : \texttt{data/interim/metadata.csv.dvc} et \texttt{data/splits/metadata\_splits.csv.dvc}.

Actuellement, l’implémentation reste partielle :
\begin{itemize}
    \item DVC est initialisé et les artefacts métadonnées sont tracés
    \item La configuration de remote n’est pas encore définie dans ".dvc/config"
    \item Aucun pipeline déclaratif n’est présent sous forme de \texttt{dvc.yaml} / \texttt{dvc.lock}
\end{itemize}

Cette configuration permet déjà de garantir une traçabilité des artefacts de la fondation data, mais elle gagnerait à être complétée par une chaîne déclarative de stages. Idéalement, le pipeline DVC serait articulé en trois étapes :
\begin{enumerate}
    \item \textbf{build\_metadata} : dépendances = annotations JSON + images pré/post, sortie = \texttt{data/interim/metadata.csv}
    \item \textbf{split\_metadata} : dépendance = \texttt{data/interim/metadata.csv}, sortie = \texttt{data/splits/metadata\_splits.csv}
    \item \textbf{check\_foundation} : dépendance = tables de métadonnées, sortie = validation de cohérence structurelle
\end{enumerate}

Une telle déclaration permettrait à \texttt{dvc repro} de rejouer uniquement les stages modifiés et d'enregistrer automatiquement les empreintes des entrées et sorties dans le fichier de verrouillage. Cela renforcerait les trois propriétés clés suivantes :
\begin{itemize}
    \item reproductibilité : reconstruction à l'identique à partir d'un commit Git et d'un checkout DVC
    \item auditabilité : traçabilité univoque de chaque exemple vers une version de code et de données brutes
    \item extensibilité : ajout ultérieur d'une source de données peut s'exprimer par un stage supplémentaire sans casser les stages existants
\end{itemize}

\subsection{MLflow Model Registry}
\begin{itemize}
    \item Enregistrement de tous les modèles expérimentaux
    \item Tagging du modèle final sélectionné
    \item Traçabilité des métriques et hyperparamètres
\end{itemize}

\subsection{Intégration FastAPI (futur)}
Point de départ pour un service de prédiction :
\begin{itemize}
    \item Endpoint POST : \texttt{/predict}
    \item Inputs : pré/post crops
    \item Outputs : label + confiance + Grad-CAM
\end{itemize}

\newpage

% Section 6 : Matrices de confusion et interprétation
\section{Résultats et interprétation}

\subsection{Confusion matrices}
Les matrices de confusion des run principales montrent :
\begin{itemize}
    \item Run B améliore significativement la classe \textit{destroyed}
    \item La classe \textit{minor-damage} reste difficile
    \item \textit{no-damage} est bien prédite
\end{itemize}

\subsection{Leçons d'ingénierie}
\begin{enumerate}
    \item Un modèle plus avancé n'est pas toujours meilleur sur les données réelles
    \item La sélection du modèle doit être basée sur les performances mesurées
    \item L'imbalance des classes affecte fortement l'optimisation
\end{enumerate}

\newpage

% Section 7 : Perspectives futures
\section{Perspectives futures}

\subsection{Court terme}
\begin{itemize}
    \item Compléter l'intégration DVC (pipeline déclaratif, remote)
    \item Déployer l'API FastAPI
    \item Rédiger la documentation technique détaillée
\end{itemize}

\subsection{Moyen/long terme}
\begin{itemize}
    \item Monitoring des prédictions en production
    \item Gestion des cas difficiles et feedback loops
    \item Extensibilité à d'autres types de désastres
    \item Auto-labeling ou semi-supervised learning
\end{itemize}

\newpage

% Section 8 : Sources de référence
\section{Références}

\subsection{Fichiers clés du projet}
\begin{itemize}
    \item \texttt{src/data/} : Data Foundation
    \item \texttt{src/training/} : Pipeline d'apprentissage
    \item \texttt{src/evaluation/} : Validation et confusion matrices
    \item \texttt{src/xai/} : Grad-CAM et explainability
    \item \texttt{src/models/} : Architecture Siamese et variantes
\end{itemize}

\subsection{Documentation interne}
\begin{itemize}
    \item \texttt{docs/training\_pipeline\_explained.md}
    \item \texttt{docs/mlflow\_registry\_serving.md}
\end{itemize}

\newpage

% Instructions pour ajouter des annexes PDF
\newpage
\appendix

\section{Annexes - PDFs du projet}

\subsection{Instructions pour intégrer les PDFs}

Pour ajouter des PDFs à ce document, procédez comme suit :

\begin{enumerate}
    \item Placez vos fichiers PDF dans le dossier \texttt{pfedocument/pdfs/}
    \item Utilisez la syntaxe suivante pour chaque PDF :
\end{enumerate}

\begin{verbatim}
\subsection{Nom du PDF}
\includepdf[pages=-]{pdfs/mon_fichier.pdf}
\end{verbatim}

\textbf{Exemple :}
\begin{verbatim}
\subsection{Rapport d'évaluation}
\includepdf[pages=-]{pdfs/evaluation_report.pdf}

\subsection{Architecture modèle}
\includepdf[pages=-]{pdfs/model_architecture.pdf}
\end{verbatim}

\vspace{1cm}

\textbf{Notes :}
\begin{itemize}
    \item \texttt{pages=-} inclut toutes les pages du PDF
    \item Utilisez \texttt{pages=1-3} pour inclure les pages 1 à 3 seulement
    \item Les PDFs doivent être en format PDF standard
\end{itemize}

\end{document}
